% GENERATED FILE. Edit canonical lessons, then run npm run generate:books.
% canonical-source-hash: fnv1a64:1729e021f0d9c11c
% canonical-lessons: ML-C283-atiyantara, ML-C283-pradhanappetta, ML-C283-suraksitamaya, ML-C283-talkkalika, ML-C283-nerittu

\chapter{Urgent, Important, Safe, Temporary, Directly}
\label{ch:ml-urgent-important-safe-temporary-directly}
\hlchaptermodality{283}

\begin{chapteropening}
\textbf{By the end of this chapter:} \emph{I can say urgent, important, safe, temporary and directly.}

Complete the last lesson of chapter 283: I can say urgent, important, safe, temporary and directly.
\end{chapteropening}

\section[\texorpdfstring{\ml{അടിയന്തര} (aṭiyantara)}{അടിയന്തര (aṭiyantara)}]{\ml{അടിയന്തര} (aṭiyantara) — urgent, emergency}

\label{lesson:ML-C283-atiyantara}



\begin{quote}
Before the new one: say the Malayalam for an attempt, then the Malayalam for success.
\end{quote}

\subsection*{You'll want to know: \ml{അടിയന്തര}}
\textbf{\ml{അടിയന്തര}} — \emph{aṭiyantara} — \textquotedblleft{}urgent, emergency\textquotedblright{}.

\begin{cousinweb}
\textbf{\ml{അടിയന്തര}} (\emph{aṭiyantara}) is of uncertain origin; neither DEDR nor CDIAL lists it.
\end{cousinweb}

\begin{grammarlens}[title={What you've built}]
One more word for this chapter.
\end{grammarlens}

\subsection*{Your turn}
\begin{itemize}
\raggedright
  \item \emph{Say it:} \emph{aṭiyantara}
  \item \emph{Say it:} \emph{aṭiyantara}, once more
  \item \emph{Say it:} say \emph{aṭiyantara}
  \item \emph{Recall:} say the Malayalam for an attempt, then the Malayalam for success, then say \emph{aṭiyantara} again
\end{itemize}

\begin{tcolorbox}[breakable,colback=teal!4,colframe=teal!35!black,title={Before you move on}]
What does \ml{അടിയന്തര} mean? (\textquotedblleft{}Urgent, emergency\textquotedblright{}.) Say it once more.
\end{tcolorbox}
\section[\texorpdfstring{\ml{പ്രധാനപ്പെട്ട} (pradhānappeṭṭa)}{പ്രധാനപ്പെട്ട (pradhānappeṭṭa)}]{\ml{പ്രധാനപ്പെട്ട} (pradhānappeṭṭa) — important}

\label{lesson:ML-C283-pradhanappetta}



\begin{quote}
Before the new one: say the Malayalam for success, then the Malayalam for urgent.
\end{quote}

\subsection*{You'll want to know: \ml{പ്രധാനപ്പെട്ട}}
\textbf{\ml{പ്രധാനപ്പെട്ട}} — \emph{pradhānappeṭṭa} — \textquotedblleft{}important\textquotedblright{}.

\begin{cousinweb}
\textbf{\ml{പ്രധാനപ്പെട്ട}} (\emph{pradhānappeṭṭa}) joins Sanskrit \emph{pradhāna}, \textquotedblleft{}chief, principal\textquotedblright{}, to native \emph{peṭṭa}, a form of the verb \textquotedblleft{}to fall into, happen\textquotedblright{} (DEDR 3853).
\end{cousinweb}

\begin{grammarlens}[title={What you've built}]
One more word for this chapter.
\end{grammarlens}

\subsection*{Your turn}
\begin{itemize}
\raggedright
  \item \emph{Say it:} \emph{pradhānappeṭṭa}
  \item \emph{Say it:} \emph{pradhānappeṭṭa}, once more
  \item \emph{Say it:} say \emph{pradhānappeṭṭa}
  \item \emph{Recall:} say the Malayalam for success, then the Malayalam for urgent, then say \emph{pradhānappeṭṭa} again
\end{itemize}

\begin{tcolorbox}[breakable,colback=teal!4,colframe=teal!35!black,title={Before you move on}]
What does \ml{പ്രധാനപ്പെട്ട} mean? (\textquotedblleft{}Important\textquotedblright{}.) Say it once more.
\end{tcolorbox}
\section[\texorpdfstring{\ml{സുരക്ഷിതമായ} (surakṣitamāya)}{സുരക്ഷിതമായ (surakṣitamāya)}]{\ml{സുരക്ഷിതമായ} (surakṣitamāya) — safe}

\label{lesson:ML-C283-suraksitamaya}



\begin{quote}
Before the new one: say the Malayalam for urgent, then the Malayalam for important.
\end{quote}

\subsection*{You'll want to know: \ml{സുരക്ഷിതമായ}}
\textbf{\ml{സുരക്ഷിതമായ}} — \emph{surakṣitamāya} — \textquotedblleft{}safe\textquotedblright{}.

\begin{cousinweb}
\textbf{\ml{സുരക്ഷിതമായ}} (\emph{surakṣitamāya}) joins Sanskrit \emph{surakṣita}, \textquotedblleft{}well protected\textquotedblright{} (CDIAL 13499), to native \emph{āya}, \textquotedblleft{}having become\textquotedblright{} (DEDR 333). DEDR 333 compares Sanskrit \emph{ām}.
\end{cousinweb}

\begin{grammarlens}[title={What you've built}]
One more word for this chapter.
\end{grammarlens}

\subsection*{Your turn}
\begin{itemize}
\raggedright
  \item \emph{Say it:} \emph{surakṣitamāya}
  \item \emph{Say it:} \emph{surakṣitamāya}, once more
  \item \emph{Say it:} say \emph{surakṣitamāya}
  \item \emph{Recall:} say the Malayalam for urgent, then the Malayalam for important, then say \emph{surakṣitamāya} again
\end{itemize}

\begin{tcolorbox}[breakable,colback=teal!4,colframe=teal!35!black,title={Before you move on}]
What does \ml{സുരക്ഷിതമായ} mean? (\textquotedblleft{}Safe\textquotedblright{}.) Say it once more.
\end{tcolorbox}
\section[\texorpdfstring{\ml{താൽക്കാലിക} (tālkkālika)}{താൽക്കാലിക (tālkkālika)}]{\ml{താൽക്കാലിക} (tālkkālika) — temporary}

\label{lesson:ML-C283-talkkalika}



\begin{quote}
Before the new one: say the Malayalam for important, then the Malayalam for safe.
\end{quote}

\subsection*{You'll want to know: \ml{താൽക്കാലിക}}
\textbf{\ml{താൽക്കാലിക}} — \emph{tālkkālika} — \textquotedblleft{}temporary\textquotedblright{}.

\begin{cousinweb}
\textbf{\ml{താൽക്കാലിക}} (\emph{tālkkālika}) $\leftarrow$ Sanskrit \emph{tātkālika}, \textquotedblleft{}lasting for that time\textquotedblright{}, taken over whole as a tatsama. Its parts are \emph{tat}, \textquotedblleft{}that\textquotedblright{}, and \emph{kāla}, \textquotedblleft{}time\textquotedblright{}.
\end{cousinweb}

\begin{grammarlens}[title={What you've built}]
One more word for this chapter.
\end{grammarlens}

\subsection*{Your turn}
\begin{itemize}
\raggedright
  \item \emph{Say it:} \emph{tālkkālika}
  \item \emph{Say it:} \emph{tālkkālika}, once more
  \item \emph{Say it:} say \emph{tālkkālika}
  \item \emph{Recall:} say the Malayalam for important, then the Malayalam for safe, then say \emph{tālkkālika} again
\end{itemize}

\begin{tcolorbox}[breakable,colback=teal!4,colframe=teal!35!black,title={Before you move on}]
What does \ml{താൽക്കാലിക} mean? (\textquotedblleft{}Temporary\textquotedblright{}.) Say it once more.
\end{tcolorbox}
\section[\texorpdfstring{\ml{നേരിട്ട്} (nēriṭṭŭ)}{നേരിട്ട് (nēriṭṭŭ)}]{\ml{നേരിട്ട്} (nēriṭṭŭ) — directly, in person}

\label{lesson:ML-C283-nerittu}



\begin{quote}
Before the new one: say the Malayalam for safe, then the Malayalam for temporary.
\end{quote}

\subsection*{You'll want to know: \ml{നേരിട്ട്}}
\textbf{\ml{നേരിട്ട്}} — \emph{nēriṭṭŭ} — \textquotedblleft{}directly, in person\textquotedblright{}.

\begin{cousinweb}
\textbf{\ml{നേരിട്ട്}} (\emph{nēriṭṭŭ}) is native Dravidian: \emph{nēr}, \textquotedblleft{}straight, direct\textquotedblright{} (DEDR 3770, with Tamil \emph{nēr}), and \emph{iṭṭu}, \textquotedblleft{}having put\textquotedblright{} (DEDR 442).
\end{cousinweb}

\begin{grammarlens}[title={What you've built}]
One more word for this chapter.
\end{grammarlens}

\subsection*{Your turn}
\begin{itemize}
\raggedright
  \item \emph{Say it:} \emph{nēriṭṭŭ}
  \item \emph{Say it:} \emph{nēriṭṭŭ}, once more
  \item \emph{Say it:} say \emph{nēriṭṭŭ}
  \item \emph{Recall:} say the Malayalam for safe, then the Malayalam for temporary, then say \emph{nēriṭṭŭ} again
\end{itemize}

\begin{tcolorbox}[breakable,colback=teal!4,colframe=teal!35!black,title={Before you move on}]
What does \ml{നേരിട്ട്} mean? (\textquotedblleft{}Directly, in person\textquotedblright{}.) Say it once more.
\end{tcolorbox}
